
\subsection{Einfache Herleitung}
Bevor wir uns an die komplexe Herleitung für elliptische Bahnen wagen, betrachten wir den vereinfachten Spezialfall.
Wir nehmen an, der Planet bewege sich auf einer perfekten Kreisbahn um die Sonne.
In diesem Fall entspricht die grosse Halbachse $a$ exakt dem konstanten Radius $r$ der Kreisbahn.

Für eine stabile Kreisbahn muss ein Kräftegleichgewicht herrschen.
Die Gravitationskraft $F_G$, die zwischen der Sonne (Masse $M$) und dem Planeten (Masse $m$) wirkt, fungiert hierbei als die benötigte Zentripetalkraft $F_Z$.
Diese Kraft zwingt den Planeten auf seine Kreisbahn.

Die Gravitationskraft nach Newton lautet:

\begin{equation}
    F_G
    = G \frac{m \cdot M}{r^2}
\end{equation}

Hierbei steht $G$ für die universelle Gravitationskonstante.
Die Zentripetalkraft für einen Körper auf einer Kreisbahn berechnet sich durch:

\begin{equation}
    F_Z
    = m \frac{v^2}{r}
    = m \omega^2 r
\end{equation}

Wir ersetzen nun die Winkelgeschwindigkeit $\omega$ durch den Ausdruck $\frac{2\pi}{T}$, wobei $T$ die Umlaufzeit des Planeten ist.
Dadurch lässt sich die Zentripetalkraft folgendermassen umschreiben:

\begin{equation}
    F_Z
    = m \frac{4\pi^2}{T^2} r
\end{equation}

Im nächsten Schritt setzen wir die Zentripetalkraft mit der Gravitationskraft gleich ($F_Z = F_G$):

\begin{equation}
    m \frac{4\pi^2}{T^2} r
    = G \frac{m \cdot M}{r^2}
\end{equation}

Die Masse des Planeten $m$ kann auf beiden Seiten der Gleichung gekürzt werden.
Anschliessend formen wir die Gleichung um.
Wir sortieren die Terme so, dass die Quadrate der Umlaufzeiten ($T^2$) und die Kuben der Radien ($r^3$) gemeinsam auf einer Seite stehen:

\begin{align}
    \frac{4\pi^2}{T^2} r
     & = G \frac{M}{r^2}
    \\
    \frac{4\pi^2}{G \cdot M}
     & = \frac{T^2}{r^3}
\end{align}

Wir erhalten als Resultat dieser Umformung:

\begin{equation}
    \frac{T^2}{r^3}
    = \frac{4\pi^2}{G \cdot M}
    = \text{konst.}
\end{equation}

Dies ist das \textbf{dritte Keplersche Gesetz} für Kreisbahnen.
Es besagt, dass der Quotient aus dem Quadrat der Umlaufzeit und der dritten Potenz des Bahnradius für alle Planeten eines Zentralgestirns eine Konstante ergibt.

Obwohl diese Herleitung mathematisch sehr verständlich ist, gilt sie nur für perfekte Kreise.
Da Planetenbahnen in der Realität Ellipsen sind, ist für einen exakten physikalischen Beweis ein allgemeinerer Ansatz notwendig.
Dieser wird im folgenden Kapitel vorgestellt.
